\documentclass[11pt]{article}

\usepackage[margin=1in]{geometry}
\usepackage{amsmath,amssymb,amsthm}
\usepackage{mathtools}
\usepackage{hyperref}

\newtheorem{proposition}{Proposition}
\newtheorem{remark}{Remark}

\DeclareMathOperator{\Spec}{Spec}

\begin{document}

\title{Equations for the Cover $M(2,2,4,8)$ over $A(2,2,4,4)$}
\author{}
\date{}
\maketitle

Let the universal $A(2,2,4,4)$-family be given by
\begin{align}
    y_1^2&=(a^2-c^2)(a^2-d^2), \label{eq:A2244-y1}\\
    y_2^2&=(b^2-c^2)(b^2-d^2), \label{eq:A2244-y2}\\
    w^2&=(a^2-c^2)(b^2-c^2). \label{eq:A2244-w}
\end{align}
Introduce
\[
    t=\frac{y_1y_2}{w}.
\]
Then
\[
    t^2=(a^2-d^2)(b^2-d^2).
\]
Equivalently, one may work on the affine model with variables
$a,b,c,d,y_1,y_2,w,t$ and equations
\begin{align}
    y_1^2&=(a^2-c^2)(a^2-d^2), \label{eq:model-y1}\\
    y_2^2&=(b^2-c^2)(b^2-d^2), \label{eq:model-y2}\\
    w^2&=(a^2-c^2)(b^2-c^2), \label{eq:model-w}\\
    wt&=y_1y_2. \label{eq:model-t}
\end{align}
Equation~\eqref{eq:model-t} implies
\[
    t^2=(a^2-d^2)(b^2-d^2)
\]
on the open locus where \eqref{eq:model-y1}, \eqref{eq:model-y2}, and
\eqref{eq:model-w} hold.

Now define
\[
    \rho=\frac ab,
    \qquad
    \sigma=\frac{a^2-c^2}{w},
    \qquad
    \tau=\frac{a^2-d^2}{t}.
\]
Then
\[
    \sigma^2=\frac{a^2-c^2}{b^2-c^2},
    \qquad
    \tau^2=\frac{a^2-d^2}{b^2-d^2}.
\]
After the usual fractional linear transformation, the associated genus two
curve is isomorphic to
\[
    Y^2=X(X+1)(X+\rho^2)(X+\sigma^2)(X+\tau^2).
\]
The chosen order $4$ point is the half of $P_0=(0,0)-\infty$ on this
normalized model.

\begin{proposition}
The cover $M(2,2,4,8)\to A(2,2,4,4)$ obtained by dividing this chosen
order $4$ point by $2$ is given, over the above open locus, by adjoining
square roots $q_1,q_2,q_3,q_4$ satisfying
\begin{align}
    q_1^2&=(1+\rho)(1+\sigma)(1+\tau), \label{eq:q1-rational}\\
    q_2^2&=\rho(1+\rho)(\rho+\sigma)(\rho+\tau), \label{eq:q2-rational}\\
    q_3^2&=\sigma(1+\sigma)(\rho+\sigma)(\sigma+\tau), \label{eq:q3-rational}\\
    q_4^2&=\tau(1+\tau)(\rho+\tau)(\sigma+\tau). \label{eq:q4-rational}
\end{align}
Equivalently,
\[
M(2,2,4,8)
=
\Spec_{A(2,2,4,4)}
\mathcal O[q_1,q_2,q_3,q_4]/(q_1^2-F_1,\ldots,q_4^2-F_4),
\]
where
\begin{align*}
F_1&=(1+\rho)(1+\sigma)(1+\tau),\\
F_2&=\rho(1+\rho)(\rho+\sigma)(\rho+\tau),\\
F_3&=\sigma(1+\sigma)(\rho+\sigma)(\sigma+\tau),\\
F_4&=\tau(1+\tau)(\rho+\tau)(\sigma+\tau).
\end{align*}
\end{proposition}

\begin{proof}
The halving criterion for the normalized curve
\[
    Y^2=X(X+1)(X+\rho^2)(X+\sigma^2)(X+\tau^2)
\]
is applied to the four square roots
\[
    1,\rho,\sigma,\tau.
\]
It says precisely that the half of the order $4$ point exists after adjoining
square roots of the four quantities
\begin{align*}
    &(1+\rho)(1+\sigma)(1+\tau),\\
    &\rho(1+\rho)(\rho+\sigma)(\rho+\tau),\\
    &\sigma(1+\sigma)(\rho+\sigma)(\sigma+\tau),\\
    &\tau(1+\tau)(\rho+\tau)(\sigma+\tau).
\end{align*}
This gives equations \eqref{eq:q1-rational}--\eqref{eq:q4-rational}.
\end{proof}

For a polynomial affine model, set
\[
    \alpha=a^2-c^2,
    \qquad
    \beta=a^2-d^2.
\]
Then
\[
    \sigma=\frac{\alpha}{w},
    \qquad
    \tau=\frac{\beta}{t}.
\]
Clearing denominators in \eqref{eq:q1-rational}--\eqref{eq:q4-rational}
gives the following equations:
\begin{align}
    bwt\,q_1^2
    &=(a+b)(w+\alpha)(t+\beta), \label{eq:q1-cleared}\\
    b^4wt\,q_2^2
    &=a(a+b)(aw+b\alpha)(at+b\beta), \label{eq:q2-cleared}\\
    bw^4t\,q_3^2
    &=\alpha(w+\alpha)(aw+b\alpha)(t\alpha+w\beta), \label{eq:q3-cleared}\\
    bwt^4\,q_4^2
    &=\beta(t+\beta)(at+b\beta)(t\alpha+w\beta). \label{eq:q4-cleared}
\end{align}
Thus an affine equation set for the desired cover is given by
\eqref{eq:model-y1}--\eqref{eq:model-t} together with
\eqref{eq:q1-cleared}--\eqref{eq:q4-cleared}.

\begin{remark}
The space $M(2,2,4,8)$ is not literally a closed subvariety of
$A(2,2,4,4)$ in the ordinary sense. Rather, it is a finite cover of
$A(2,2,4,4)$ obtained by adjoining the square roots $q_1,q_2,q_3,q_4$.
Generically this is a $(\mathbb Z/2)^4$-cover.
\end{remark}

Multiplying the four square-root equations gives
\[
q_1^2q_2^2q_3^2q_4^2
=
\rho\sigma\tau
\left(
(1+\rho)(1+\sigma)(1+\tau)
(\rho+\sigma)(\rho+\tau)(\sigma+\tau)
\right)^2.
\]
Therefore the first natural intermediate double cover is
\[
    r^2=\rho\sigma\tau.
\]
In the original variables, this is equivalent, up to multiplication by a
square, to
\[
    r^2=abwt=ab\,y_1y_2.
\]
Thus the necessary condition noticed from the product relation is
\[
    abwt\in \mathbb Q^{*2},
\]
or, since $wt=y_1y_2$,
\[
    ab\,y_1y_2\in \mathbb Q^{*2}.
\]
This intermediate double cover is the first visible step in the tower of
double covers from $M(2,2,4,8)$ down to $A(2,2,4,4)$.

\end{document}
